\chapter{Introduction}
\label{ch:introduction}

A user interface (UI) transition graph represents application screens and the transitions caused
by user actions. Each edge records a triggering event, a guard and a destination. The graph supports
reachability queries and helps identify screens affected by a component change.

Transition graphs support Android test generation and exploration
\cite{yang2015wtg,goalexplorer2019,storydroid2019} and crawling-based web testing
\cite{crawljax2012,feedex2013}. Their links to source code and tests also support change impact
analysis \cite{arnold1993impact,bohner1996software}.

\section{Motivation}

Dynamic extraction executes an application and records the states and transitions it observes.
Its coverage depends on test data, authentication and the exploration budget, and each new version
requires another run. Static extraction derives a graph from source code without executing the
application. It can run on each commit, but recovers only the constructs its analysis supports and
must approximate targets it cannot resolve.

These approaches are complementary. Much of the navigation in the applications studied here is
behind authentication guards. A static graph can identify candidate screens, transitions and
guards for a crawler, while execution can confirm observed transitions. This thesis evaluates only
static extraction; the combination is left for future work (Section~\ref{sec:future}).

Static extraction is established for Android. Window transition graph analyses derive windows and
transitions from bytecode and link callback control flow to the widgets registering each handler
\cite{yang2015wtg,yang2015callbacks,yang2018wtg}. Early web navigation models combined templates
and configuration with crawling or server execution \cite{riccatonella2001,webuml2004}; a patent
describes a fully static method \cite{tkachuk2013patent}.

Modern single-page applications have mostly been modelled dynamically. Crawljax builds a state
flow graph by firing events against the running Document Object Model (DOM) \cite{crawljax2012}.
FeedEx and FRAGGEN improve exploration while retaining execution as the source of the model
\cite{feedex2013,fraggen2022}. A survey covering 2014--2025 reviews model-based web testing
\cite{websurvey2025}; every listed tool whose model construction was examined for this thesis
executes the application. Framework-specific static analyses remain few. One constructs a React
component dependency graph without routing \cite{kurant2025}. Three extract navigation:
\begin{itemize}
\item Oshima and Kishi extract Angular page transitions for model checking. Their stated
restrictions exclude nested routes and route parameters, and their model omits page transitions
caused by data-binding changes \cite{oshima2018model}.
\item Rima et al.\ extract Angular workflows as a coverage denominator for generated end-to-end
tests \cite{rima2026e2e}.
\item Zhang et al.\ analyse compiled single-page application bundles to check access-control
policies against reachable screens \cite{zhang2025veriflow}.
\end{itemize}
Section~\ref{sec:dynamic-spa} compares these approaches.

Early Ajax applications lacked declarative route tables, so crawlers used DOM states to observe
navigation. Angular, React Router and Vue Router now expose route tables, template navigation and
service calls in source code; routes can also name guards \cite{angulardocs}. These declarations
provide inputs for static extraction. Existing studies evaluate extraction against hand-built
references on their chosen applications, but do not measure losses construct by construct.

Navigation graphs also motivate applications beyond model-based testing. GPTDroid uses a
question-and-answer loop over the current mobile screen and reports higher activity coverage and
more detected bugs than its best baseline on 93 applications \cite{liu2024make}. AutoDroid supplies
a model with an application-specific transition graph built by offline dynamic exploration
\cite{wen2024autodroid}; LLMDroid uses a model to guide exploration towards uncovered behaviour
\cite{llmdroid2025}. Related work explores web interfaces using vision-language models
\cite{wang2024leveraging}, and a survey covers large language models in software testing
\cite{wang2024software}. Coding agents may also use navigation structure: SWE-bench tasks require
coordinated edits across repository files \cite{jimenez2024swebench}, and repository interfaces
affect agent performance \cite{yang2024sweagent}. Whether navigation graphs benefit these consumers
is not tested here.

This thesis builds an Angular extractor to measure recovery and classify losses. It proposes
\emph{scope attribution}, a dataflow analysis over screen sets, for the largest loss category under
the study's rule order. The analysis is specified independently of the extractor. Measurement is
the main contribution; scope attribution is a modest extension of earlier mechanisms. Its added
step places navigation in services, guards and effects, recovering \AblationFullAddedSites{} sites
(\AblationFullAddedShare{}) above outlet attribution. This step combines a member-level service
call graph, a fallback join and stand-ins for guards and effects. The call graph alone adds
\AblationCallGraphNetSites{} sites (\AblationCallGraphNetShare{}) over a per-class join
(Section~\ref{sec:res-ablation}).

The extractor supports nested, parameterised and lazily loaded route trees as well as scope
attribution. The stated restrictions of the earlier Angular model-checking approach exclude the
first two (Section~\ref{sec:res-earlier}). The extractor was developed and evaluated on the same
corpus, so its measurements describe these applications rather than estimate performance on
Angular applications generally. Section~\ref{app:repairs} records the development corrections;
Section~\ref{sec:limitations} discusses their implications. Chapter~\ref{ch:discussion} separates
implementation gaps from limits of the approach.

\section{Research questions and scope}
\label{sec:rqs}

The aim is to measure how far static extraction of a UI transition graph gets on real Angular
applications, why it fails where it does, and what scope attribution adds. Five research questions
(RQs) follow from it; the measures are defined in Section~\ref{sec:metrics} and answered in
Chapter~\ref{ch:results}.

\begin{description}[nosep,leftmargin=2.6em,labelwidth=2.2em]
\item[RQ1] How much of the navigation in Angular applications can static analysis recover as
  deterministic transitions?
\item[RQ2] Where does static extraction fail, and why?
\item[RQ3] How does scope attribution affect the recovered share and precision of static
  navigation analysis?
\item[RQ4] Does the extracted graph respond to source changes in the constructs it represents, as
  change impact analysis requires?
\item[RQ5] Are the losses a property of the applications or of the tool?
\end{description}

The questions are answered on six open-source Angular applications pinned to fixed commits
(Chapter~\ref{ch:evaluation}), with the final version of the extractor. The thesis does not claim
completeness: it states a soundness property for the must tier that holds only relative to stated assumptions (Claim~\ref{claim:must}) and lists what it approximates
(Section~\ref{sec:soundiness}). It does not measure the graph's effect on any consumer, it evaluates
Angular only, and its figures on the corpus are in-sample, since the extractor was developed on it.

\section{Thesis structure}

Figure~\ref{fig:overview} summarises the study.

\begin{figure}[t]
\centering
\begin{tikzpicture}[node distance=5mm and 7mm]
  % input
  \node[box, minimum width=27mm, fill=mutedlight] (src) {Angular source\\[-1pt]{\footnotesize routes, templates,}\\[-3pt]{\footnotesize classes}};
  % extractor
  \node[box, right=12mm of src, minimum width=34mm] (cons) {construct discovery};
  \node[box, above=3mm of cons, minimum width=34mm] (routes) {route extraction};
  \node[box, below=3mm of cons, minimum width=34mm, hl] (attr) {scope attribution\\[-2pt]{\footnotesize (Chapter~\ref{ch:analysis})}};
  \begin{scope}[on background layer]
    \node[draw=muted, rounded corners=4pt, fit=(routes)(cons)(attr), inner sep=5pt,
      label={[font=\footnotesize\itshape, text=muted]above:uigraph extractor}] (ext) {};
  \end{scope}
  \node[box, right=12mm of cons, minimum width=28mm] (graph) {UI transition graph\\[-2pt]{\footnotesize must / may / unknown}\\[-3pt]{\footnotesize + witnesses}};
  \draw[->, thick] (src.east) -- (src.east -| ext.west);
  \draw[->] (routes) -- (cons);
  \draw[->] (cons) -- (attr);
  \draw[->, thick] (ext.east |- graph) -- (graph.west);

  % evaluation row, well below the extractor
  \coordinate (row) at ($(ext.south)+(0,-17mm)$);
  \node[box, minimum width=26mm, font=\footnotesize, anchor=north] (l) at (row) {loss breakdown\\[-2pt]defect classes};
  \node[box, minimum width=26mm, font=\footnotesize, left=2mm of l] (c) {route parser\\[-2pt]node coverage};
  \node[box, minimum width=26mm, font=\footnotesize, left=2mm of c] (b) {witness check\\[-2pt]must-tier precision};
  \node[box, minimum width=26mm, font=\footnotesize, right=2mm of l] (f) {mutation\\[-2pt]detection rate};
  \node[box, minimum width=16mm, font=\footnotesize, right=2mm of f] (d) {E2E\\[-2pt]oracle};
  \node[font=\footnotesize\itshape, text=muted, anchor=south west] at ($(b.north west)+(0,8mm)$) {evaluation (Chapter~\ref{ch:evaluation})};
  % bus from the graph down to the evaluation row
  \coordinate (bus) at ($(row)+(0,6mm)$);
  \draw[muted] (graph.south) |- (bus -| d.north);
  \draw[muted] (bus -| d.north) -- (bus -| b.north);
  \foreach \n in {b,c,l,f,d} \draw[->, muted] (bus -| \n.north) -- (\n.north);
  % research questions
  \node[below=1mm of b, font=\scriptsize] {RQ1, RQ3};
  \node[below=1mm of c, font=\scriptsize] {RQ1};
  \node[below=1mm of l, font=\scriptsize] {RQ2, RQ3};
  \node[below=1mm of f, font=\scriptsize] {RQ4};
  \node[below=1mm of d, font=\scriptsize] {RQ1};
\end{tikzpicture}
\caption{Overview of the study. The extractor turns Angular source into a UI transition
graph; scope attribution is the analysis this thesis adds. Five measurements, each tied to
the research questions it answers, evaluate the graph.}
\label{fig:overview}
\end{figure}


Chapter~\ref{ch:background} reviews transition graphs, related extraction methods, Angular and
the analysis foundations. Chapter~\ref{ch:analysis} defines the model, its conditional soundness
claim, scope attribution and the implementation. Chapter~\ref{ch:evaluation} describes the corpus,
metrics, evaluation methods and protocol corrections. Chapter~\ref{ch:results} answers the research
questions and reports statically detected dead-link candidates. Chapter~\ref{ch:discussion}
interprets the results and discusses limitations, framework transfer and future work.
Chapter~\ref{ch:conclusion} summarises the contributions.
